\scanpage{63}
\mspage{12}

\uline{Lemme} Soit $(X,x_0)$ dans $\mathrm{Cat}_{\bullet}$. Si $X$ est $W$-asphérique, alors $\Omega(X,x_0)\overset{\text{déf}}{=}\underline{\mathrm{Ch}}(X;x_0,x_0)$ \deleted{est} l'est aussi (moyennant les axiomes $\textcircled{1}\textcircled{2}\textcircled{3}\textcircled{4}$ sur $W$). Donc (par itération) les $\Omega^n(X,x_0)$ sont $W$-asphériques $(\forall n\in\mathbb N)$.

% slanted
\margin{Utilise $W(1),W(2~\mathrm{bis})$ $W(3),W(4\,\mathrm b'')$ sur $W$}

\uline{Corollaire} Si de plus $W$ satisfait l'axiome $\textcircled{5}$ de connexité, alors $X$ est \deleted{\unsure{donc}} $W_\infty$-asphérique.

En effet, ``on sait'' que cela signifie que tous les $\Omega^n(X,x_0)$ sont $0$-connexes. Or l'axiome de connexité $\textcircled{5}$ assure que $W$-asphérique implique $0$-connexe.

% slanted
\margin{C'est prouvé fort simplement dans le corollaire au ``\unsure{th, 1|th. 1}'', page 190}

Mais le ``on sait'' est un tout petit peu léger$^*$. Pour le justifier, il faudrait, ou bien que je fasse le lien avec la théorie \unsure{$\frac{1}{2}$ simpliciale|ksimpliciale|simpliciale} classique, et prouve que les $\Omega^n$ que je construis sont ``les bons'', i.e. que leur Nerf redonne bien (à isom près dans $\mathrm{Hot}$) les $\Omega^n$ classiques (??). En particulier, que le $\pi_0(\Omega^n(X,x))$ est bien le $\pi_n(X,x_0)$ ``de tout le monde''. Ou bien que je prouve directement que la $0$-connexité des \unsure{\S} $\Omega^n(X,x_0)$ équivaut à la $W_\infty$-asphéricité de $X$, définie comment ? \unsure{grossièrement} : le Artin-Mazur.

C'est donc un travail qu'il me faudra faire \unsure{absolument|obscurément}, Mais ce n'est plus qu'un simple travail de routine, et je \deleted{\unsure{puis}}\inserted{\unsure{peux}} énoncer dès maintenant le

% margin note written sideways
\margin{les hypothèses utiles sur $W$ sont finalement $W(1),W(2~\mathrm{bis}),W(3),W(4\,\mathrm b''),$ et $W\ne\mathrm{Fl}(\mathrm{Cat})$}

% slanted
\margin{Il suffit même de $\textcircled{4}$ sous une forme affaiblie, une des formulations ``faibles'' \unsure{des nos|de mes} \unsure{deux} ``\unsure{absolus|absol.}'' cf le NB.}
% slanted
\margin{\unsure{ci|ce} suffit, cf. \unsure{infra} p. 13, 14.}

\uline{Théorème 1}. Soit $W\subset\mathrm{Fl}(\mathrm{Cat})$ satisfaisant les conditions $\textcircled{1}\textcircled{2}\textcircled{3}\textcircled{4}\textcircled{5}$ (axiomes de saturation, de localisation, de l'objet final, de fibration, de connexité). Alors $W\subset W_\infty$.

\margin{\uline{NB} Au lieu de $W(5)$, il suffit de supposer $W\ne\mathrm{Fl}(\mathrm{Cat})$, cf \deleted{\unsure{th. 2}} p. 20.}

\uline{Corollaire} On peut caractériser $W_\infty$ comme la plus grande partie de $\mathrm{Fl}(\mathrm{Cat})$ satisfaisant aux

% archive sheet number: 63
